% !TeX spellcheck = hu_HU
% !TeX encoding = UTF-8
% !TeX program = xelatex
%----------------------------------------------------------------------------
\section{Programozható statikus analízis eszköz}\label{sec:ftr:blast}
%----------------------------------------------------------------------------

Az egyik alapvető motiváció a nyelv megvalósítás mögött egy olyan eszköz alapjának létrehozása volt, amelyik lehetővé teszi a felhasználója számára, hogy saját maga által írt kódrészleteket tudjon futtatni egy C vagy C++ projektben való különböző absztrakt nyelvi elemek megtalálása esetén.
Ennek eleme, hogy ezeknek a nyelveknek az \gls{AST}-jében való kereséshez egy kényelmesen használható \gls{DSL}-t lehessen megvalósítani, de a találat esetén futó kód ne legyen limitálva ezen \gls{DSL} elvárásai szerint.

Példának egy AWK jellegű eszközt lehetséges elképzelni.
Ez azonban egy alacsonyabb absztrakciós szinten -- a szimpla fájlok primitív szöveges tartalma felett -- biztosít futtatható kódrészleteket, amikor egy-egy reguláris kifejezést sikeresen megtalál a bemenetén.
Ennek vitathatatlan előnye, hogy a lehető legáltalánosabb és így gyakorlatilag minden szöveges fájlon működőképes.
Azonban emiatt nem rendelkezik megfelelő mennyiségű szemantikus információval az esetlegesen (szemi-)strukturált bemenetéről, így csak a legáltalánosabb egyszerű szöveg alapján tud keresni.

Ez olyan esetekben, ahol a szövegnek szigorúan meghatározott formátuma van kiemelkedően zavaró tud lenni, főleg ha esetlegesen szóköz és egyéb tagoló karakterekre viszont érzéketlen a szöveg formátuma.
Ilyen például a \gls{GPL}-ek nagy része (\pl C, C++, Java, C4, stb.) de azon kevés nyelv (programozási vagy egyéb), amelyik legalább részlegesen tagolásra érzékeny (\pl Python, Haskell) is ugyanebben a problémába ütközik.
Kiemelkedő probléma, hogy olyan szemi-strukturált adattároló formátumok, mint az egyes \gls{XML} dialektusok, és a \gls{HTML} sem kezelhető megfelelően általános módon reguláris kifejezésekkel még elméletileg sem; egyszerűen ezen nyelvek nemregularitása miatt.

Az előző paragrafus által felvázolt problémának az alapja, hogy a különféle nyelvekkel leírt programokat (vagy szemi-strukturált adatokat) a használt nyelv konkrét szintaxisával leírva tároljuk a forrásfájlokban.
Emiatt sok extra szintaktikai elemet -- ami az emberi felhasználóknak akár még hasznos is -- kell tudni megfelelő módon kezelni, amitől a feldolgozás csak nehezebbé változik.
Ebben rejlik egy nagy erőssége a már többször említett LISP nyelveknek, mert az általuk használt \gls{sexpr} formulák gyakorlatilag nem tartalmaznak semmilyen olyan extra szintaktikát, ami a feldolgozáshoz nem direkt módon tartozik.

A nyelvek nagy többségének azonban a konkrét szintaxisa nem ilyen egyszerűen kezelhető. 
A C és C++ nyelvek például szintaktikailag kifejezetten nehezen kezelhetőek.
Egyik példa az egyszerű szöveges feldolgozás által tapasztalt nehézségekre, hogy nem különböztetik meg a globális és egy függvényen belüli lokális változókat, sőt C++-ban még az osztályok tagváltozói is sok esetben azonos formátumúak.
Emiatt, ha azt az egyszerű feladatot szeretnénk megvalósítani, hogy megkeressük az összes globális változót egy adott projekten belül, akkor nem fogjuk tudni extra szemantikus információ nélkül.
Összefoglalva, ha valamilyen szemantikai alapú leírással szeretnénk megkülönböztetni elemeket valamilyen nyelvben megadott szövegben, akkor az AWK-nál erősebb eszközre van szükség, amelyik képes az absztrakt szintaxisban végrehajtani a keresést.

A statikus analízis eszközök emiatt sokszor maguk is ilyen eszközökre építenek.
Valamilyen könyvtár segítéségével feldolgozzák a konkrét szintaxist, és az így kapott \gls{AST}-ben, mint gráfban, keresnek olyan mintázatokat, amelyekről azt gondolják, hogy nem megfelelő viselkedést okoznak futás közben.
Egy C és C++ nyelveket támogató nyílt forráskódú eszköz a {\tt clang-tidy} az LLVM projekt gondozásában, de léteznek kommersz eszközök is.

Sokszor azonban egy zárt, előre lefejlesztett, és a binárisba -- vagy szolgáltatásba -- csomagolt ellenőrzés listát támogatnak ezek a megvalósítások.
A felhasználói oldalon ezeket csak ki- és bekapcsolni tudjuk, esetlegesen pár ellenőrzés által használt konstansokat van lehetőségünk felülírni valamilyen projekt szintű konfigurációs fájlban.

Létezhetnek azonban olyan statikusan ellenőrizhető elvárások, amelyek egy-egy projekt sajátosságai, és így a fejlesztőinek és csak nekik, hasznos lenne, ha ezeket az elvárásokat lehetséges lenne automatikusan ellenőriztetni.
Három példát ismertetek erre.

Egy dinamikus könyvtárat valósít meg egy projekt C vagy C++ nyelven.
Ilyenkor a dinamikusan betöltődő könyvtár (\gls{SO} vagy \gls{DLL}), kell hogy rendelkezzen egy publikusan hívható függvényekből álló interfésszel.
Ezeket ilyenkor kiemelkedően körültekintő módon kell elkészíteni: a megfelelő hívási konvencióval, konzisztens nevekkel, és kikapcsolt névátalakítással (mangling) C++-ban.
Ennek egyik része lehet, hogy a projekt rövidítésével kell, hogy kezdődjön, minden publikus függvény neve.
Utóbbira nem fognak a projekt fejlesztői tudni egyből elérhető ellenőrzést találni, hiszen a projektjük neve egyedi.
Így valahogyan nekik szükséges megvalósítani a megfelelő ellenőrzést.

A másik példa egy nagyobb projekt esetén érvényes, amelyiknek fontos, hogy platform független módon, operációs rendszerek között forduljon és fusson.
Emiatt rendelkeznek egy saját absztrakciós réteggel a számukra szükséges funkcionalitással, és ennek különböző platformokra írt megvalósításokkal.
Ilyenkor fontos lehet ellenőrizni, hogy az absztrakciós réteg által nyújtott funkcionalitást nem kerüli-e ki véletlenül az üzleti logikát megvalósító réteg.
Ehhez ismerni kell a projekt által fejlesztett absztrakciós réteget, illetve a projekt struktúrát, ami alapján lehetségessé válik eldönteni, hogy mely kódrészletek tartoznak az absztrakciós rétegbe, és melyikek az üzleti logikába.
Emiatt ebben az esetben is saját kezűleg kell a projekt fejlesztőinek megvalósítani az ellenőrzéseiket.

Az utolsó egy még jobban szakterület-specifikus működést feltételez: ez oktatási környezetben jelenik meg.
Tegyük fel, hogy valamilyen tárgy keretein belül szükség van a diákoknak/hallgatóknak egy vagy több házi feladatot elkészíteni C vagy C++ programozási nyelven, a tárgy által megadott extra követelményeket és elvárásokat betartva.
Ez lehet akár még feladatonként eltérő is: például egy programozást bevezető tárgyban a rekurziót gyakoroltató háziban nem megengedett a különböző ciklusok használata.
Bár erre lehetséges szöveges ellenőrzést megvalósítani ezen nyelvek esetén ezt könnyű kijátszani, \ld \aref{lst:cpp-tricks}.~listázás.
Emiatt emberi felülvizsgálat lehet szükséges.
Egyszerűbb esetekben is -- a csalni próbáló megoldókat nem számolva -- hasznos lehet az egyes tárgykövetelmények automatikus ellenőrzése.
Egy ilyen jellegű példa lehet, ha a tárgy biztosít egy memóriaszivárgás ellenőrző eszközt, amelyiken keresztül kell a dinamikus memóriát kezelni. Ilyenkor lehet keresni olyan allokációs hívásokat, amelyek nem ezen keresztül történnek.

\begin{lstlisting}[gobble=2,numbers=left,label=lst:cpp-tricks,caption={Szöveges ellenőrzést kijátszó ,,for'' ciklus},language=c]
	#include <stdio.h>
	#define PRIMITVE_CAT(a, b) a##b
	
	int main() {
		PRIMITVE_CAT(fo,r) (int i = 0; i < 10; ++i)
		    printf("i: %d\n", i);
	}
\end{lstlisting}

Az utolsó példánál kiemelkedően látszik, hogy létezhetnek olyan ellenőrzések, amelyek egy projekt számára elengedhetetlenül hasznosak, mindazonáltal más projekteben semmi értelme nem lenne futtatni azokat: ,,Egy komplex üzleti logikával rendelkező, C-ben írt, több platformot támogató natív grafikus ablakban futó rendszer, amelyik fejlesztésében nem megengedettek a ciklusok.''

Ezekben az esetekben, érthető módon, nem lehetséges a kész termékként elérhető eszközökbe belefejleszteni minden egyes projekt saját elvárásait, mert egy kezelhetetlenül nagy ellenőrzés halmaz keletkezne belőle, úgy hogy az így elérhető ellenőrzések nagy részét csak az az egy-két projekt használná, amelyik eredetileg kérvényezte annak lefejlesztését.
Ez kiemelkedően igaz az ingyenes és nyílt fejlesztésű megoldásokra, hiszen nem lesz sok ember, aki ezt számunkra karban fogja tartani, ha ezt nem tesszük meg magunk.
Összességében ebben a helyzetben is azok az ellenőrzések lennének nagyobb mértékben használva, amelyek a konkrét projekt elvárásoktól függetlenül biztosítanak hasznos ellenőrzéseket: nem bezárt fájl, biztosan elszivárgott memória, tautologikus elágazás feltételek, stb.

Még mielőtt az általam eltervezett C4 alapú megoldást bemutatom megemlítendő, hogy egyes eszközök esetében lehetséges saját ellenőrzésekkel kibővíteni az ellenőrző eszközt.

Ilyen az előzőekben már említett {\tt clang-tidy} program, amelyik egy C++-ba beépített \gls{eDSL}-en keresztül teszi lehetővé, hogy ellenőrzéseket lehessen definiálni.
A folyamat elég részletesen dokumentált \cite{llvmprojectGettingInvolvedExtra2025}.
A C++-ban megvalósított \gls{eDSL} az Clang \gls{AST}-n való keresést biztosító \gls{AST}Matchers interfész, ami rendelkezik saját dokumentációval \cite{llvmprojectASTMatcherReference2025}.
Mivel ez az \gls{eDSL} C++ osztályokat és függvényhívás operátorokat használ, így ha egyes képességekkel nem rendelkezik, a felhasználó számára lehetséges bővíteni.
Amint beépítésre került az új ellenőrzés, a {\tt clang-tidy} programot újra kell fordítani.
Innentől szabadon használható, a feladat csupán a folyamatosan fejlődő LLVM és Clang verziókkal lépést tartva karban tartani.

Egy másik opcióként az új verziói a {\tt clang-tidy} eszköznek lehetővé teszik az előbb említett \gls{eDSL} kódot használni, úgy hogy azt egy szöveges mezőbe a konfigurációs fájlba beírjuk, majd definiáljuk a megfelelő kulcs-érték párokat a generálni kívánt diagnosztikai üzenetekhez \cite{llvmprojectQueryBasedCustom2025}.
Ez a megközelítés jóval dinamikusabb, mint az előbb említett, de jelentősen gyengébb is: az \gls{eDSL} már nem rendelkezik a C++ által biztosított kifejezőerővel, és így komplexebb ellenőrzéseket (\pl egy sztring tartalmának része-e egy adott rész-sztring) nem lehetséges megvalósítani benne.
A fejlesztői élmény is rosszabb, hiszen a kódot egy konfigurációs fájlban egy szövegben írjuk, egy tényleges forrásfájl helyett.
Egyéb nehézség, bár ez tranziens, hogy ez a lehetőség jelenleg kísérleti állapotban van a következő LLVM/Clang verzió forrásában, így a fejlesztők jelenleg nem garantálnak kompatibilitást az elkövetkező verziókkal.

Ezen problémák kiküszöbölésére tervezek megoldás adni, egy olyan C4-et használó \gls{DSL} alapú eszközt fejlesztésével, amelyik a különböző projekt szintű elvárások ellenőrzését teszi lehetővé.
Fontos elvárás, hogy az ellenőrzések megvalósításához ne legyen szükség egy közös statikus analízis eszköz forráskódjába beépíteni, és ne konfigurációs fájlokba rejtett szöveges mezőkbe kelljen kódot írni.
\Aref{lst:blast-ex}.~listázásban bemutatom a jelenlegi tervek szerinti \gls{DSL} használatát, amelyik a C4 gazdanyelvben kerül megvalósításra.

\begin{lstlisting}[gobble=4,numbers=left,label=lst:blast-ex,caption={Példa lehetséges \gls{DSL}-re egy statikus analízis eszköz számára}]
    let my_checks checks
        match variable definition where { |def| is_global def }
            do { |def|
                println strcat "globális változó: "
                                name_of def
            }

        match function call named "malloc" where { |call| file_of call != "leak_tracer.h" }
            do { |call|
                println strfmt "nem felügyelt malloc: {}:{}" file_of call, line_of call
            }
    done
    perform my_checks in "projekt/mappa"
\end{lstlisting}




